% =====================================================================
% Homework 03 --- Analise do perfil com XFOIL
% Incluido por main.tex. Tabelas em tex_analise_xfoil/ e figuras em
% resultados_analise_xfoil/ (mesmo nivel que este arquivo e o main.tex).
% Codigos em analise_xfoil/.
% =====================================================================

\section{Analise Aerodinamica com XFOIL}
\label{sec:xfoil}

\subsection{Objetivo}

Descreva o que foi calculado com o XFOIL: polar aerodinamica,
curva $C_L(\alpha)$, arrasto, momento, comparacao entre perfis ou qualquer
outra analise pedida no laboratorio.

\subsection{Configuracao}

Registre aqui as condicoes de escoamento e a forma como a analise foi
conduzida. Exemplos:
\begin{itemize}
\item numero de Reynolds;
\item numero de Mach;
\item faixa de angulo de ataque;
\item criterio de convergencia;
\item configuracao de transicao, se aplicavel.
\end{itemize}

\begin{table}[H]
\centering
\caption{Parametros da analise com o XFOIL.}
\label{tab:xfoil-config}
\begin{tabular}{ll}
\toprule
Parametro & Valor \\
\midrule
Reynolds & preencher \\
Mach & preencher \\
Faixa de $\alpha$ & preencher \\
Observacoes & preencher \\
\bottomrule
\end{tabular}
\end{table}

\subsection{Resultados}

Apresente aqui os graficos e tabelas obtidos nas corridas do XFOIL.

\begin{figure}[H]
\centering
% \includegraphics[width=0.82\textwidth]{resultados_analise_xfoil/polar.png}
\caption{Substitua esta legenda por uma descricao do principal resultado
do XFOIL.}
\label{fig:xfoil-polar}
\end{figure}

\subsection{Discussao}

Comente o comportamento do perfil, a faixa linear de sustentacao,
o arrasto, o ponto de perda, a sensibilidade a Reynolds/Mach e quaisquer
diferencas em relacao a modelos mais simples ou a resultados esperados.
